\clearpage\chapter{Arakelov theory}
\section*{1. Overview}
\subsection*{Intuitive}
\paragraph{The problem}

Number theory treats every prime number as a finite place, but a number field also has ordinary real or complex size. Arakelov theory was developed to put these two kinds of information into one geometric framework. It adds a place at infinity to arithmetic geometry, equips the complex points with metrics, and defines intersections, heights, and Riemann--Roch formulas that combine finite prime data with analytic data.



\subsection*{Formal}
Let $K$ be a number field with ring of integers $\mathcal O_K$. Arakelov geometry treats finite prime ideals and archimedean embeddings as places of the same arithmetic space. An arithmetic divisor combines a finite divisor with real-valued data at the complex places, usually represented by a Green function or Hermitian metric. Its degree and intersection number are sums of finite algebraic contributions and archimedean analytic contributions.
\section*{2. History}
\subsection*{Historical development}
\begin{historylist}
\paragraph{History and motivation}

Suren Arakelov introduced the theory in the 1970s to study Diophantine geometry in higher dimensions. The motivating analogy is between a smooth projective curve over a finite field and an arithmetic surface over the integers. Prime ideals describe finite fibers, but $\operatorname{Spec}(\mathbb Z)$ has no prime ideal representing the ordinary absolute value. Arakelov's idea was to compactify the arithmetic picture by adding an infinite fiber and replacing missing algebraic information with a Hermitian metric on the complex points.

This was not an attempt to decorate number theory with analysis for its own sake. The goal was to make theorems for number fields resemble corresponding theorems for function fields. Intersections, degrees, and positivity should see every place of the field, including the real and complex ones \cite{arakelov_theory_ref2,arakelov_theory_ref4}.

\end{historylist}
\section*{3. Theory}
\subsection*{Core theory}
\paragraph{Finite places and the place at infinity}

Let $K$ be a number field and $\mathcal O_K$ its ring of integers. A nonzero prime ideal $\mathfrak p$ gives a finite place and a valuation measuring divisibility by $\mathfrak p$. The field also has embeddings into $\mathbb R$ or $\mathbb C$, called archimedean places. For $\mathbb Q$, the finite places come from primes $p$, while the infinite place is the usual absolute value $|\cdot|$.

For a nonzero rational number $x$, the product formula says that the product of its normalized absolute values over all places is 1:
\[
 \prod_v |x|_v=1.
\]
In logarithmic form, the sum of the local logarithms is zero. This identity is the arithmetic reason that a global height can ignore a common rescaling of homogeneous coordinates. It also illustrates the theory's central principle: a quantity that looks incomplete at one place becomes balanced after all places are included.

\paragraph{Arithmetic surfaces and metrised bundles}

An arithmetic surface is, roughly, a regular model $\mathfrak X$ of a smooth projective curve over $K$ defined over $\operatorname{Spec}(\mathcal O_K)$. Its finite fibers lie over prime ideals. After an embedding $K\to\mathbb C$, the complex points form a Riemann surface $X(\mathbb C)$. Arakelov theory equips line bundles or vector bundles on this complex surface with Hermitian metrics, compatible with complex conjugation.

The metric is the analytic replacement for the missing point at infinity. A metrised line bundle remembers both an algebraic divisor at finite places and a smooth metric at archimedean places. Its arithmetic degree measures finite multiplicities and the logarithmic size contributed at infinity. Green functions describe the potential created by a divisor and allow intersections involving the infinite component to be defined.

\paragraph{A simple height example}

If $p/q$ is a rational number in lowest terms, its logarithmic height is
\[
 h(p/q)=\log\max(|p|,|q|).
\]
The numerator and denominator record finite-prime information, while the maximum is an archimedean size. If one writes the same projective point as $(cp,cq)$, the product formula explains why the normalized global height does not depend on the chosen scaling. For algebraic numbers, the height averages the logarithmic absolute values over all embeddings and places.

\paragraph{Divisors with an infinite part}

For a field $K$, an arithmetic surface $\mathfrak X$, and a chosen complex embedding, an arithmetic or $c$-divisor has the form
\[
 D=\sum_i k_i C_i+\sum_{\sigma}\lambda_\sigma X_\sigma.
\]
Here the $C_i$ are irreducible codimension-one closed subsets of the arithmetic model, $k_i$ are integer coefficients at finite places, and $X_\sigma$ denotes the complex or real component associated with an archimedean embedding, with real coefficient $\lambda_\sigma$. The sum respects conjugate pairs of complex embeddings. These divisors form an abelian group.

The ordinary divisor records where a rational function vanishes or has poles. The infinite coefficients record the metric or Green-function contribution. Principal divisors satisfy a global compatibility relation because the product formula prevents an isolated local contribution from being arbitrary.

\paragraph{Arithmetic intersection theory}

Arakelov defined an intersection pairing for divisors on arithmetic surfaces attached to smooth projective curves over number fields. At finite places, two divisors intersect in the usual algebraic way on a fiber. At infinity, their Green functions contribute an analytic integral. The sum is a global arithmetic intersection number.

This lets one formulate an arithmetic Hodge index theorem, self-intersection inequalities, and height pairings. Jean-Benoît Bost developed potential-theoretic methods using Green functions and obtained arithmetic Lefschetz-type results. Gerd Faltings established an arithmetic Riemann--Roch theorem, a Noether formula, a Hodge index theorem, and positivity results for the dualizing sheaf. These are the arithmetic counterparts of familiar results for algebraic surfaces \cite{arakelov_theory_ref3,arakelov_theory_ref5}.

\paragraph{Arithmetic Chow groups}

An arithmetic cycle of codimension $p$ is a pair $(Z,g)$: an algebraic $p$-cycle $Z$ on the arithmetic scheme and a Green current $g$ that describes its archimedean contribution. Green currents generalize Green functions and can have controlled logarithmic singularities along $Z$.

The arithmetic Chow group $\widehat{\mathrm{CH}}^p(\mathfrak X)$ is obtained by quotienting arithmetic cycles by arithmetic versions of rational equivalence and trivial cycles. It is the natural home for arithmetic intersection products and Chern classes. The notation is heavier than the basic idea: algebraic cycles and analytic corrections are identified whenever they represent the same global arithmetic object.

\paragraph{Arithmetic Riemann--Roch}

Ordinary Grothendieck--Riemann--Roch describes how the Chern character changes under a proper map. In schematic form,
\[
 \operatorname{ch}(f_*E)=f_*(\operatorname{ch}(E)\operatorname{Td}(T_{X/Y})).
\]
The arithmetic version uses metrised bundles and arithmetic characteristic classes:
\[
 \widehat{\operatorname{ch}}(f_*[E])
 =f_*\bigl(\widehat{\operatorname{ch}}(E)
 \widehat{\operatorname{Td}}^{\,R}(T_{X/Y})\bigr).
\]
The modified Todd class contains an additional characteristic series involving zeta values and derivatives. The correction is the analytic price of including the infinite places. In applications, the theorem relates arithmetic degrees, determinants of cohomology, and heights of bundles under a smooth proper map \cite{arakelov_theory_ref6}.

\section*{4. Important theorems and results}
\begin{enumerate}[leftmargin=*,itemsep=0.35em]

\item \textbf{Product formula.} The product of normalized absolute values of a nonzero number over all places of a number field is one. It makes global heights invariant under scaling.

\item \textbf{Arakelov intersection theory.} Divisors on arithmetic surfaces have a global intersection number containing both finite algebraic terms and archimedean Green-function terms.

\item \textbf{Arithmetic Hodge index theorem.} Under suitable hypotheses, the arithmetic intersection pairing satisfies a negativity or positivity statement analogous to the classical Hodge index theorem.

\item \textbf{Arithmetic Riemann--Roch.} Pushforward of a metrised bundle is governed by arithmetic Chern and Todd classes, with an additional zeta-function correction.

\item \textbf{Faltings' finiteness theorem.} The arithmetic-geometric methods surrounding Arakelov theory contribute to finiteness theorems for objects such as abelian varieties over a number field with prescribed dimension and bounded bad reduction.
\end{enumerate}
\noindent\emph{Source for these results:} \cite{arakelov_theory_ref1}.

\section*{5. Applications}
Arakelov theory combines finite-prime information with archimedean geometry. Heights and arithmetic intersection numbers then measure the global size and interaction of arithmetic objects.
\begin{enumerate}[leftmargin=*,itemsep=0.35em]
\item \textbf{Diophantine equations.} Heights measure the global complexity of rational and algebraic points. A bound on height turns the search for infinitely many possible points into a finite search, and height inequalities can show that only finitely many points satisfy a Diophantine condition.
\item \textbf{Elliptic and abelian varieties.} Canonical heights combine local contributions and behave quadratically under multiplication. They help distinguish torsion points from points of infinite order and are central in arithmetic geometry.
\item \textbf{Arithmetic surfaces.} Intersections of sections, fibers, and metrised divisors quantify how arithmetic curves meet. The resulting inequalities are tools for proving finiteness and positivity statements.
\item \textbf{Analytic number theory.} Archimedean metrics and Green functions connect arithmetic intersection numbers with complex analysis, theta functions, and determinants of Laplacians.
\item \textbf{Generalizations.} Modern developments include Arakelov intersection theory in higher dimensions, Hodge--Arakelov theory, $p$-adic Hodge theory, and attempts such as $\mathbb F_1$ geometry to model a geometry over a hypothetical field with one element.
\item \textbf{Cryptographic height functions.} In lattice-based cryptography, the geometric height of an algebraic number measures the size of its minimal polynomial coefficients, which directly governs the hardness of lattice problems such as the shortest-vector problem. Arakelov-theoretic heights combine archimedean embeddings (Euclidean norm on coefficient vectors) with non-archimedean valuations at primes (divisibility of coefficients), yielding a single global complexity measure. Cryptographers use height bounds to estimate the security parameter of ring-based schemes such as NTRU and ring-LWE: the discriminant of the number field controls lattice volume, while the Mahler measure of the defining polynomial, an archimedean height quantity, controls the distribution of lattice points. Arakelov intersection pairings on arithmetic surfaces attached to these rings further predict how splitting behavior at finite places affects the algebraic structure available for key generation, providing a principled way to select number fields that resist the fastest known attacks \cite{arakelov_theory_ref3,arakelov_theory_ref4}.
\end{enumerate}


\theoryconnections{arakelov_theory}{%
\item Arakelov theory combines the finite-prime information of algebraic number theory with the ordinary complex geometry at the infinite places.
\item A divisor records zeros and poles algebraically, a line bundle packages that divisor, and a metric on the complex points records its archimedean size.
\item Intersection theory can then pair a finite contribution with a Green-function contribution, producing an arithmetic intersection number.
\item Heights arise from exactly this pairing, so they measure both divisibility at primes and size in the complex plane \cite{arakelov_theory_ref3,arakelov_theory_ref4}.
}{%
\item The resulting heights let Diophantine geometry compare solutions of an equation by a single global size.
\item On an elliptic curve, height inequalities can show that only finitely many rational points occur or can bound the search for them.
\item Arithmetic Riemann--Roch and arithmetic Chow groups use the same finite-plus-infinite bookkeeping to transfer geometric theorems into arithmetic statements, while arithmetic dynamics measures the growth of a point under repeated iteration \cite{arakelov_theory_ref1,arakelov_theory_ref5,arakelov_theory_ref6}.
}

\section*{Glossary}
\begin{description}[leftmargin=*,style=nextline,itemsep=0.3em]

\item[\textbf{Arithmetic surface.}] A regular model of a smooth projective curve defined over the ring of integers of a number field, whose finite fibers lie over prime ideals and whose complex points form a Riemann surface.

\item[\textbf{Place.}] Either a finite prime ideal of a number field or an archimedean embedding into $\mathbb R$ or $\mathbb C$; Arakelov theory treats both kinds on an equal footing.

\item[\textbf{Green function.}] A real-valued function attached to a divisor that describes its analytic behavior at archimedean places, serving as the metric counterpart to the algebraic data at finite primes.

\item[\textbf{Height.}] A global measure of the arithmetic size of a point or divisor, combining finite-prime valuations with archimedean absolute values; heights bound the search for rational solutions to Diophantine equations.

\item[\textbf{Product formula.}] The identity that the product of all normalized absolute values of a nonzero element over every place of a number field equals one, reflecting a perfect balance between finite and infinite places.

\item[\textbf{Arithmetic Chow group.}] A group of arithmetic cycles---pairs of algebraic cycles and Green currents---modulo an arithmetic equivalence relation, serving as the natural home for arithmetic intersection products.

\item[\textbf{Valuation.}] A function measuring the divisibility of elements of a number field by a prime ideal; it is the local building block of arithmetic.

\item[\textbf{Hermitian metric.}] A smooth inner product on a complex vector bundle that varies continuously over the base; in Arakelov theory it provides archimedean data.

\item[\textbf{Line bundle.}] A vector bundle of rank one over a variety or scheme; a divisor determines a line bundle.

\item[\textbf{Archimedean place.}] An embedding of a number field into R or C, measuring the ordinary absolute value; the infinite counterpart to a finite prime ideal.

\end{description}

\section*{7. Sample Questions}
\emph{Exercise reference:} \cite{arakelov_theory_ref1}. The cited edition is the source for the exercise set; consult its Exercises section for the official statement and numbering.
\begin{enumerate}[leftmargin=*,itemsep=0.9em]

\item \textbf{Exercise 1.} Verify the product formula for $x=6$ over $\mathbb{Q}$, i.e., compute $|6|_v$ for every place $v$ and confirm the product equals 1.

\emph{Solution.} The archimedean place gives $|6|_\infty=6$. At $p=2$: $|6|_2=|2^1\cdot 3|_2=2^{-1}=1/2$. At $p=3$: $|6|_3=3^{-1}=1/3$. At all other primes $|6|_p=1$. The product is $6\cdot(1/2)\cdot(1/3)\cdot 1\cdots=1$.

\item \textbf{Exercise 2.} Compute the logarithmic Weil height $h(p/q)$ for the rational numbers $5/2$, $7/1$, and $1/12$.

\emph{Solution.} By definition $h(p/q)=\log\max(|p|,|q|)$ for $p/q$ in lowest terms. Thus $h(5/2)=\log 5$, $h(7/1)=\log 7$, and $h(1/12)=\log 12$.

\item \textbf{Exercise 3.} Let $K=\mathbb{Q}(i)$. List all places of $K$ and state how many archimedean places $K$ has.

\emph{Solution.} The places of $K$ correspond to prime ideals of $\mathbb{Z}[i]$ (the finite places) and embeddings $K\hookrightarrow\mathbb{C}$ up to complex conjugation (the archimedean places). Since $K$ has one complex embedding pair, there is $r_1=0$ real places and $r_2=1$ complex place, giving exactly $1$ archimedean place. The finite places are the prime ideals $(1+i)$, $(3)$, $(2+i)$, $(2-i)$, and so on for each rational prime.

\item \textbf{Exercise 4.} Compute the height $h(\alpha)$ for $\alpha=1+\sqrt{5}$ viewed as an element of $\mathbb{Q}(\sqrt{5})$, using both conjugate embeddings $\sigma_1: \sqrt{5}\mapsto\sqrt{5}$ and $\sigma_2:\sqrt{5}\mapsto-\sqrt{5}$.

\emph{Solution.} The Weil height of an algebraic number is $h(\alpha)=\frac{1}{[K:\mathbb{Q}]}\sum_v n_v\log\max(|\sigma_v(\alpha)|,1)$ summed over all places. The two conjugate values are $1+\sqrt{5}\approx 3.236$ and $1-\sqrt{5}\approx -1.236$. Each archimedean place contributes $\log(1+\sqrt{5})+\log|1-\sqrt{5}|=\log(1+\sqrt{5})+\log(\sqrt{5}-1)=\log((1+\sqrt{5})(\sqrt{5}-1))=\log(5-1)=\log 4$. Dividing by $[K:\mathbb{Q}]=2$: $h(\alpha)=\frac{1}{2}\cdot 2\log 2=\log 2$.

\item \textbf{Exercise 5.} The degree of a divisor on an arithmetic surface includes both finite and archimedean contributions. If a divisor $D$ has finite part with degree $3$ at all non-archimedean places and archimedean contribution $-1$, what is the arithmetic degree $\widehat{\deg}(D)$?

\emph{Solution.} The arithmetic degree is the sum of the finite algebraic degree and the archimedean correction. Here $\widehat{\deg}(D)=3+(-1)=2$.

\end{enumerate}
\section*{8. References}
\begin{thebibliography}{99}
\bibitem{arakelov_theory_ref1} J. H. Silverman, \emph{The Arithmetic of Elliptic Curves}.
\bibitem{arakelov_theory_ref2} J. Neukirch, \emph{Algebraic Number Theory}.
\bibitem{arakelov_theory_ref3} S. Lang, \emph{Introduction to Arakelov Theory}.
\bibitem{arakelov_theory_ref4} E. Bombieri and W. Gubler, \emph{Heights in Diophantine Geometry}.
\bibitem{arakelov_theory_ref5} G. Faltings, \emph{Calculus on Arithmetic Surfaces}.
\bibitem{arakelov_theory_ref6} H. Gillet and C. Soulé, "An arithmetic Riemann--Roch theorem," \emph{Inventiones Mathematicae}.
\end{thebibliography}
